% ============================================================
%  PACKAGES & CONFIGURATION — IACS Template
%  Intelligence Artificielle et Cybersécurité
%  ENSA Béni Mellal
% ============================================================

% --- Encodage et langue ---
\usepackage[utf8]{inputenc}
\usepackage[T1]{fontenc}
\usepackage[french]{babel}
\usepackage{csquotes}

% --- Polices ---
\usepackage{mathptmx}          % Times New Roman (corps du texte)
\usepackage{helvet}            % Helvetica pour les titres
\usepackage{courier}           % Courier pour le code

% --- Géométrie de la page ---
\usepackage[
  a4paper,
  top=2.5cm,
  bottom=2.5cm,
  left=2.5cm,
  right=2.5cm,
  headheight=15pt
]{geometry}

% --- Couleurs IACS ---
\usepackage[table,xcdraw,dvipsnames]{xcolor}

% Palette principale IACS — Bleu foncé #000099
\definecolor{iacsPrimary}{HTML}{000099}      % Bleu foncé principal
\definecolor{iacsSecondary}{HTML}{000066}    % Bleu très foncé
\definecolor{iacsAccent}{HTML}{3366CC}       % Bleu moyen (accents, liens)
\definecolor{iacsTeal}{HTML}{000099}         % Bleu principal (puces, titres)
\definecolor{iacsLight}{HTML}{E8EEFF}        % Bleu très clair (fonds)
\definecolor{iacsGreen}{HTML}{06D6A0}        % Vert IA
\definecolor{iacsOrange}{HTML}{F77F00}       % Orange alerte (cyber)
\definecolor{iacsRed}{HTML}{D62828}          % Rouge alerte
\definecolor{iacsGray}{HTML}{6C757D}         % Gris neutre
\definecolor{iacsLightGray}{HTML}{E9ECEF}    % Gris clair (fonds tableaux)
\definecolor{iacsDark}{HTML}{212529}         % Quasi-noir (texte)
\definecolor{iacsTableHead}{HTML}{000099}    % En-tête tableaux
\definecolor{iacsChapterBg}{HTML}{E8EEFF}    % Fond page de chapitre

% --- Graphiques et figures ---
\usepackage{graphicx}
\usepackage{float}
\usepackage{subcaption}
\usepackage{wrapfig}
\usepackage{tikz}
\usetikzlibrary{shapes.geometric, arrows.meta, positioning, calc, shadows, decorations.markings}

% --- Tableaux ---
\usepackage{booktabs}
\usepackage{multirow}
\usepackage{multicol}
\usepackage{longtable}
\usepackage{array}
\usepackage{colortbl}
\usepackage{tabularx}

% --- Listes ---
\usepackage{enumitem}
\setlist[itemize]{
  label=\textcolor{iacsTeal}{\textbullet},
  leftmargin=1.5em,
  itemsep=3pt
}
\setlist[itemize,2]{
  label=\textcolor{iacsAccent}{--},
  leftmargin=1.5em
}
\setlist[enumerate]{
  label=\textcolor{iacsTeal}{\arabic*.},
  leftmargin=1.5em,
  itemsep=3pt
}

% --- Code source ---
\usepackage{listings}
\lstdefinestyle{iacscode}{
  backgroundcolor=\color{iacsPrimary!5},
  basicstyle=\ttfamily\small\color{iacsDark},
  keywordstyle=\bfseries\color{iacsTeal},
  stringstyle=\color{iacsGreen},
  commentstyle=\itshape\color{iacsGray},
  numberstyle=\tiny\color{iacsGray},
  numbers=left,
  numbersep=8pt,
  frame=l,
  framerule=2pt,
  rulecolor=\color{iacsAccent},
  breaklines=true,
  captionpos=b,
  tabsize=2,
  showstringspaces=false
}
\lstset{style=iacscode}

% --- Algorithmes ---
\usepackage[ruled,vlined,french]{algorithm2e}
\SetAlgoLined
\SetKwInput{KwEntree}{Entrée}
\SetKwInput{KwSortie}{Sortie}

% --- Mathématiques ---
\usepackage{amsmath,amssymb,amsfonts}

% --- En-têtes et pieds de page ---
\usepackage{fancyhdr}
\pagestyle{fancy}
\fancyhf{}
\fancyhead[L]{\small\textcolor{iacsGray}{\leftmark}}
\fancyhead[R]{\small\textcolor{iacsGray}{PFA -- IACS 2025/2026}}
\fancyfoot[C]{\thepage}
\fancyfoot[R]{\small\textcolor{iacsGray}{ENSA Béni Mellal}}
\renewcommand{\headrulewidth}{0.5pt}
\renewcommand{\headrule}{\hbox to\headwidth{\color{iacsTeal}\leaders\hrule height \headrulewidth\hfill}}
\renewcommand{\footrulewidth}{0.3pt}
\renewcommand{\footrule}{\hbox to\headwidth{\color{iacsAccent}\leaders\hrule height \footrulewidth\hfill}}

% --- Titres des chapitres et sections ---
\usepackage{titlesec}

% Style Chapitre — bande colorée latérale
\titleformat{\chapter}[display]
  {\normalfont\Huge\bfseries\sffamily\color{iacsPrimary}}
  {\tikz[remember picture,overlay]{
    \fill[iacsTeal] (current page.north west) rectangle ([xshift=1.2cm]current page.south west);
    \node[rotate=90, anchor=south, font=\fontsize{60}{72}\selectfont\bfseries\sffamily,
          text=white, inner sep=10pt]
      at ([xshift=0.6cm, yshift=-5cm]current page.north west)
      {\thechapter};
  }}
  {0pt}
  {\filright\fontsize{28}{34}\selectfont}
  [\vspace{0.5cm}{\color{iacsAccent}\titlerule[2pt]}]

\titleformat{\section}
  {\normalfont\Large\bfseries\sffamily\color{iacsSecondary}}
  {\textcolor{iacsTeal}{\thesection}}{0.8em}{}
  [{\color{iacsAccent}\titlerule[0.8pt]}]

\titleformat{\subsection}
  {\normalfont\large\bfseries\sffamily\color{iacsSecondary}}
  {\textcolor{iacsTeal}{\thesubsection}}{0.8em}{}

\titleformat{\subsubsection}
  {\normalfont\normalsize\bfseries\color{iacsDark}}
  {\textcolor{iacsAccent}{\thesubsubsection}}{0.8em}{}

\titlespacing*{\chapter}{0pt}{-30pt}{30pt}
\titlespacing*{\section}{0pt}{18pt}{6pt}
\titlespacing*{\subsection}{0pt}{12pt}{4pt}
\titlespacing*{\subsubsection}{0pt}{8pt}{4pt}

% --- Légendes figures/tableaux ---
\usepackage[
  font=small,
  labelfont={bf,color=iacsTeal},
  textfont={color=iacsDark},
  format=hang,
  margin=1cm
]{caption}

% --- Hyperliens ---
\usepackage[
  colorlinks=true,
  linkcolor=iacsTeal,
  citecolor=iacsTeal,
  urlcolor=iacsAccent,
  bookmarks=true,
  bookmarksnumbered=true,
  pdfstartview=FitH
]{hyperref}

% --- Bibliographie ---
\usepackage[
  backend=biber,
  style=ieee,
  sorting=nyt,
  maxnames=3,
  minnames=1
]{biblatex}

% Les identifiants de fichiers Google Drive (Partie C du dossier de sources) sont de
% longues chaînes alphanumériques opaques, sans "/" interne où couper : sans ce
% réglage, biblatex ne les coupe qu'aux points/slashes et déborde largement de la
% marge dans la bibliographie. On abaisse les pénalités de coupure aux changements
% de casse/chiffre pour forcer une coupure plus fréquente sur ces URL précises.
\setcounter{biburllcpenalty}{100}
\setcounter{biburlucpenalty}{100}
\setcounter{biburlnumpenalty}{100}

% --- Table des matières ---
\usepackage{tocloft}
\renewcommand{\cftchapfont}{\bfseries\sffamily\color{iacsTeal}}
\renewcommand{\cftchappagefont}{\bfseries\color{iacsTeal}}
\renewcommand{\cftsecfont}{\sffamily\color{iacsSecondary}}
\renewcommand{\cftsecpagefont}{\color{iacsSecondary}}
\renewcommand{\cftsubsecfont}{\color{iacsDark}}
\renewcommand{\cftchapleader}{\cftdotfill{\cftdotsep}}
\setlength{\cftbeforechapskip}{8pt}

% --- Glossaire et acronymes ---
\usepackage[acronym,toc]{glossaries}
\makeglossaries

% --- Divers ---
\usepackage{lipsum}            % Texte de remplissage
\usepackage{pdfpages}          % Inclusion de PDF
\usepackage{setspace}          % Interligne
\onehalfspacing                % Interligne 1.5
\setlength{\parindent}{1cm}    % Retrait première ligne
\usepackage{etoolbox}          % Utilitaires

% --- Commandes personnalisées ---
\newcommand{\iacsnote}[1]{%
  \begin{tcolorbox}[colback=iacsLight, colframe=iacsTeal, 
    title={\textcolor{white}{\faInfoCircle\; Note}},
    coltitle=white, fonttitle=\bfseries\sffamily,
    colbacktitle=iacsTeal, boxrule=0.5pt, arc=3pt]
    #1
  \end{tcolorbox}
}

\newcommand{\iacswarning}[1]{%
  \begin{tcolorbox}[colback=iacsOrange!8, colframe=iacsOrange,
    title={\textcolor{white}{\faExclamationTriangle\; Attention}},
    coltitle=white, fonttitle=\bfseries\sffamily,
    colbacktitle=iacsOrange, boxrule=0.5pt, arc=3pt]
    #1
  \end{tcolorbox}
}

% --- Boîtes colorées ---
\usepackage[most]{tcolorbox}
\usepackage{fontawesome5}

% Redéfinir les commandes après chargement de tcolorbox
\renewcommand{\iacsnote}[1]{%
  \begin{tcolorbox}[colback=iacsLight, colframe=iacsTeal,
    title={\textcolor{white}{\faInfoCircle\; Note}},
    coltitle=white, fonttitle=\bfseries\sffamily,
    colbacktitle=iacsTeal, boxrule=0.5pt, arc=3pt]
    #1
  \end{tcolorbox}
}

\renewcommand{\iacswarning}[1]{%
  \begin{tcolorbox}[colback=iacsOrange!8, colframe=iacsOrange,
    title={\textcolor{white}{\faExclamationTriangle\; Attention}},
    coltitle=white, fonttitle=\bfseries\sffamily,
    colbacktitle=iacsOrange, boxrule=0.5pt, arc=3pt]
    #1
  \end{tcolorbox}
}

\newcommand{\iacsdefinition}[2]{%
  \begin{tcolorbox}[colback=iacsGreen!5, colframe=iacsGreen,
    title={\textcolor{white}{\faBook\; #1}},
    coltitle=white, fonttitle=\bfseries\sffamily,
    colbacktitle=iacsGreen!80!black, boxrule=0.5pt, arc=3pt]
    #2
  \end{tcolorbox}
}
